\documentclass[11pt]{article}

\usepackage[a4paper,margin=29mm]{geometry}
\usepackage{amsmath,amssymb,amsthm,mathtools}
\usepackage{array,booktabs}
\usepackage{microtype}
\usepackage[numbers,sort&compress]{natbib}
\usepackage[hidelinks]{hyperref}

\newtheorem{theorem}{Theorem}
\newtheorem{proposition}[theorem]{Proposition}
\newtheorem{corollary}[theorem]{Corollary}
\newtheorem{lemma}[theorem]{Lemma}
\newtheorem{claim}[theorem]{Claim}
\theoremstyle{definition}
\newtheorem{remark}[theorem]{Remark}
\newtheorem{question}[theorem]{Question}

\newcommand{\R}{\mathbb{R}}
\newcommand{\disc}{\Delta}
\newcommand{\eps}{\varepsilon}

\title{Machine-Dependent Chairman Assignment on Two Machines:\\
the Liu--Reis Bound Holds and Is Optimal}
\author{Angus Muffatti}
\date{24 November 2026 --- draft}

\begin{document}
\maketitle

\begin{abstract}
Liu and Reis conjectured that every fractional assignment of $n$ jobs to
$m$ machines with machine-dependent weights $d_{ij}$ can be rounded to an
integral assignment in which every machine, at every prefix of the job
sequence, deviates from its fractional load by at most
$D=\max_{i,j}d_{ij}$.  The conjecture is false in general: there is a
counterexample with eleven machines.  We prove that it is true for two
machines, and that the constant is best possible there.  The upper bound
is proved by an invariant in the style of the trapping families of Lund
and Magazinov: a parallelogram of reachable states that has the origin in
its convex hull.  It yields a linear-time algorithm, and the rounding
respects the support of the fractional assignment.  For the lower bound
we give an explicit family of two-machine instances with $2k+1$ jobs whose
optimal value is exactly $(1-2^{-(k+1)})D$; with machine-independent
weights the two-machine optimum is $D/2$.  We also show that no online
rule has a bounded guarantee, already for two machines, and that the
upper bound fails if the two machines may be charged with weights of
opposite sign.  The least number of machines for which the conjecture
fails is therefore between $3$ and $11$.
\end{abstract}

\section{Introduction}

Let $x\in[0,1]^{m\times n}$ be a fractional assignment, $\sum_i x_{ij}=1$
for every column $j$, and let $d\in\R_{\ge0}^{m\times n}$ be a matrix of
machine-dependent weights with $D=\max_{i,j}d_{ij}$.  For an integral
assignment $y\in\{0,1\}^{m\times n}$ with the same column sums, the
\emph{row-prefix discrepancy} is
\[
 \disc_i(t)=\sum_{j=1}^{t}d_{ij}\,(x_{ij}-y_{ij}),\qquad
 i\in[m],\ t\in[n].
\]
Write
\[
 \kappa_m=\sup_{\text{instances with $m$ rows}}\ \min_{y}\ \max_{i,t}\
 \frac{|\disc_i(t)|}{D},\qquad \kappa=\sup_m\kappa_m .
\]
When the weights do not depend on the machine ($d_{ij}=d_j$), Liu and
Reis \citep{LiuReis2026} proved that the minimum is always strictly below
$D$, which generalizes Tijdeman's theorem on the chairman assignment
problem \citep{Tijdeman1973}.  For machine-dependent weights they
conjectured $\kappa\le1$ (Conjecture~3 of the arXiv version of
\citep{LiuReis2026}); the case $t=n$ alone follows from the rounding
theorem of Lenstra, Shmoys and Tardos \citep{LenstraShmoysTardos1990}.
The conjecture was refuted in \citep{Muffatti2026} by an instance with
$11$ rows and $15$ columns, and \citep{MuffattiChains2026} proves
$\kappa\ge1+W(1/e)\approx1.2785$, where $W$ is the Lambert function.
These constructions need a number of rows that grows with the target
value.  This paper determines what happens with two rows.

\begin{theorem}\label{thm:upper}
Let $m=2$.  For all weights $d_{ij}\in[0,D]$ and every fractional
assignment $x$ there is an integral assignment $y$ with
\[
 |\disc_i(t)|\le D\qquad\text{for }i=1,2\text{ and all }t\in[n].
\]
Moreover $y$ can be chosen with $y_{ij}=0$ whenever $x_{ij}=0$, and it can
be computed with $O(n)$ arithmetic operations.
\end{theorem}

\begin{theorem}\label{thm:lower}
For every $k\ge0$ there is an instance with two rows, $2k+1$ columns, all
weights in $(0,1]$ and all fractional entries equal to $1/2$, for which
\[
 \min_y\max_{i,t}|\disc_i(t)|\;=\;1-2^{-(k+1)} .
\]
\end{theorem}

\begin{corollary}
$\kappa_2=1$.  The conjecture of Liu and Reis holds for two machines, and
no constant smaller than $D$ is valid there.
\end{corollary}

For comparison, with machine-independent weights and two rows the two
discrepancies are negatives of each other, the problem is
one-dimensional, and the exact constant is $1/2$ (this is Tijdeman's
value $1-1/(2m-2)$ at $m=2$).  Machine-dependent weights double it.

\begin{theorem}\label{thm:online}
Let $m=2$.  In the online problem, where $y_{\cdot j}$ must be fixed when
column $j$ arrives and the adversary chooses column $j+1$ after it sees
$y_{\cdot j}$, the adversary can force $\max_{i,t}|\disc_i(t)|>cD$ for
every constant $c$.
\end{theorem}

Theorem~\ref{thm:online} explains why the proof of
Theorem~\ref{thm:upper} cannot follow one state: there is no bounded
region that a single state can be kept in.  The proof instead keeps a set
of at most four reachable states.  The corresponding algorithm reads the
columns once, but it carries up to four partial assignments and chooses
between them only at the end.

\paragraph{Relation to earlier work.}
The proof of Theorem~\ref{thm:upper} adapts the \emph{trapping families}
of Lund and Magazinov \citep{LundMagazinov2015}.  They prove that in every
normed plane a sequence of vectors of norm at most $1$ admits signs for
which all partial sums have norm at most $2$.  In the notation of
Section~\ref{sec:planar} this is the case of Theorem~\ref{thm:upper} in
which every fractional entry equals $1/2$, and it holds there for vectors
of arbitrary direction.  Their trapping sets are centrally symmetric
parallelograms.  For general fractional entries the parallelogram is no
longer symmetric about the origin, and two things change: a vertex of the
parallelogram can leave the feasible square, and the directions of the
vectors matter.  Indeed the statement becomes false if the two weights of
a column may have opposite signs (Proposition~\ref{prop:mixed}).  A
general reduction from arbitrary to equal fractional entries, in the
manner of Lov\'asz, Spencer and Vesztergombi
\citep{LovaszSpencerVesztergombi1986}, loses a factor of $2$ and gives
only $2D$.

For general $m$ the selection form of the theorem of B\'ar\'any and
Grinberg \citep{BaranyGrinberg1981,Barany2008} gives $\kappa_m\le 2m$: for
each column the $m$ possible increments of the discrepancy vector have
sup norm at most $D$ and contain the origin in their convex hull.  A
bound independent of $m$ is closely related to the prefix Beck--Fiala
conjecture of Bansal, Rohwedder and Svensson
\citep{BansalRohwedderSvensson2022}, which is open.  A different
conjecture of \citep{LiuReis2026}, with machine-independent weights and a
support constraint, was refuted in \citep{BercziLiuReisTarnawski2026}.

\section{The planar form}\label{sec:planar}

Throughout Sections~\ref{sec:planar}--\ref{sec:online} we have $m=2$ and,
after scaling, $D=1$.  For column $j$ put
\[
 h_j=(d_{1j},d_{2j})\in[0,1]^2,\qquad w_j=x_{1j}\in[0,1],
\]
and let $y_j=y_{1j}\in\{0,1\}$, so that $y_j=1$ means that column $j$ is
assigned to row~$1$.  Since $x_{2j}=1-w_j$ and $y_{2j}=1-y_j$,
\[
 p_t:=\bigl(-\disc_1(t),\,\disc_2(t)\bigr)=\sum_{j=1}^{t}(y_j-w_j)\,h_j .
\]
Thus the state is a point of the plane that starts at the origin, and
column $j$ moves it by one of
\[
 +(1-w_j)\,h_j\qquad\text{or}\qquad -w_j\,h_j .
\]
The two possible end points differ by $h_j$, a vector with non-negative
coordinates and sup norm at most $1$.  Let $Q=[-1,1]^2$.  Then
$\max_{i,t}|\disc_i(t)|\le1$ if and only if $p_t\in Q$ for all $t$.  Let
$R_t$ be the set of points $p_t$ that are reachable by a path with
$p_1,\dots,p_t\in Q$; so $R_0=\{0\}$, and Theorem~\ref{thm:upper} states
that $R_n\ne\emptyset$.

\section{The upper bound}\label{sec:upper}

\paragraph{The invariant.}
For $t\ge0$ let $I(t)$ be the statement: \emph{there are
$g_1,g_2\in[0,1]^2$ and $\lambda_1,\lambda_2\in[0,1]$ such that, with
\[
 \pi(\eps)=(\eps_1-\lambda_1)\,g_1+(\eps_2-\lambda_2)\,g_2,
 \qquad\eps\in\{0,1\}^2,
\]
every one of the four points $\pi(\eps)$ that lies in $Q$ belongs to
$R_t$.}

The four points are the vertices of a parallelogram that has the origin
in its convex hull.  A vertex outside $Q$ carries no information.
$I(0)$ holds with $g_1=g_2=0$.

\begin{lemma}\label{lem:nonempty}
$I(t)$ implies $R_t\ne\emptyset$.
\end{lemma}
\begin{proof}
Each coordinate of $\pi(1,0)=(1-\lambda_1)g_1-\lambda_2g_2$ is a
difference of two numbers in $[0,1]$.  So $\pi(1,0)\in Q$, and therefore
$\pi(1,0)\in R_t$.
\end{proof}
The same argument shows $\pi(0,1)\in Q$.

\paragraph{The step.}
Assume $I(t)$, and let $g_3=h_{t+1}$ and $\lambda_3=w_{t+1}$.  For
$\eps\in\{0,1\}^3$ put
\[
 P(\eps)=\sum_{i=1}^{3}(\eps_i-\lambda_i)\,g_i
 =\pi(\eps_1,\eps_2)+(\eps_3-\lambda_3)\,g_3 .
\]
If $\pi(\eps_1,\eps_2)\in R_t$ and $P(\eps)\in Q$, then
$P(\eps)\in R_{t+1}$: extend the path by $y_{t+1}=\eps_3$.  Call $\eps$
\emph{bad} if $P(\eps)\in Q$ but $\pi(\eps_1,\eps_2)\notin Q$.  If $\eps$
is not bad and $P(\eps)\in Q$, then $\pi(\eps_1,\eps_2)\in Q$, hence
$\pi(\eps_1,\eps_2)\in R_t$ by $I(t)$, hence $P(\eps)\in R_{t+1}$.

\begin{claim}\label{cl:bad}
A bad vertex is $(0,0,1)$ or $(1,1,0)$, and at most one of these two is
bad.
\end{claim}
\begin{proof}
The points $\pi(1,0)$ and $\pi(0,1)$ lie in $Q$, so no vertex
$(1,0,\cdot)$ or $(0,1,\cdot)$ is bad.  The point
$\pi(0,0)=-\lambda_1g_1-\lambda_2g_2$ has non-positive coordinates.  If it
is outside $Q$, one of its coordinates is below $-1$, and the same
coordinate of $P(0,0,0)=\pi(0,0)-\lambda_3g_3$ is below $-1$; so
$(0,0,0)$ is not bad.  In the same way $(1,1,1)$ is not bad.  Finally,
$\pi(0,0)\notin Q$ means that a coordinate of
$\lambda_1g_1+\lambda_2g_2$ exceeds $1$, which gives
$\lambda_1+\lambda_2>1$; and $\pi(1,1)\notin Q$ gives
$(1-\lambda_1)+(1-\lambda_2)>1$.  The two cannot both hold.
\end{proof}

\begin{lemma}\label{lem:middle}
Among three vectors in the closed positive quadrant of $\R^2$, one is a
non-negative combination of the other two.
\end{lemma}
\begin{proof}
A zero vector is such a combination.  If all three are non-zero, order
them by angle.  Either all three are parallel, or the middle one lies in
the cone spanned by the other two.
\end{proof}

Fix names $\{\alpha,\beta,\gamma\}=\{1,2,3\}$ and numbers
$\rho,\sigma\ge0$ with
\begin{equation}\label{eq:rel}
 g_\beta=\rho\,g_\alpha+\sigma\,g_\gamma .
\end{equation}
The relation is symmetric in $(\alpha,\rho)$ and $(\gamma,\sigma)$.

\begin{lemma}\label{lem:slide}
\textup{(i)} There is $\tau\ge0$ such that the triple $\mu$ with
\[
 \mu_\alpha=\lambda_\alpha-\tau\rho,\qquad
 \mu_\beta=\lambda_\beta+\tau,\qquad
 \mu_\gamma=\lambda_\gamma-\tau\sigma
\]
lies in $[0,1]^3$ and satisfies: $\mu_\beta=1$, or $\rho>0$ and
$\mu_\alpha=0$, or $\sigma>0$ and $\mu_\gamma=0$.

\textup{(ii)} There is $\tau\ge0$ such that the triple $\mu$ with
\[
 \mu_\alpha=\lambda_\alpha+\tau\rho,\qquad
 \mu_\beta=\lambda_\beta-\tau,\qquad
 \mu_\gamma=\lambda_\gamma+\tau\sigma
\]
lies in $[0,1]^3$ and satisfies: $\mu_\beta=0$, or $\rho>0$ and
$\mu_\alpha=1$, or $\sigma>0$ and $\mu_\gamma=1$.

In both cases $\sum_i\mu_ig_i=\sum_i\lambda_ig_i$.
\end{lemma}
\begin{proof}
Take the largest $\tau$ for which the three numbers stay in $[0,1]$.  It
is finite because of the coordinate $\beta$, and the constraint that stops
it is one of the three that are listed.  The sum does not change, by
\eqref{eq:rel}.
\end{proof}

Suppose that $\mu\in[0,1]^3$, $\sum_i\mu_ig_i=\sum_i\lambda_ig_i$ and
$\mu_k=v\in\{0,1\}$ for some $k$.  Then
\begin{equation}\label{eq:hull}
 0=\sum_{i=1}^{3}(\mu_i-\lambda_i)\,g_i
 =\sum_{\eps:\ \eps_k=v}\Bigl(\prod_{i\ne k}\mu_i^{\eps_i}(1-\mu_i)^{1-\eps_i}\Bigr)P(\eps),
\end{equation}
so the origin is a convex combination of the four points $P(\eps)$ of the
face $\{\eps_k=v\}$ of the cube.  Hence each outcome of
Lemma~\ref{lem:slide} supplies such a face: $\mu_\beta=1$ gives
$\{\eps_\beta=1\}$, $\mu_\alpha=0$ gives $\{\eps_\alpha=0\}$, and so on.

\begin{claim}\label{cl:face}
There are $k\in\{1,2,3\}$, $v\in\{0,1\}$ and $\mu\in[0,1]^3$ with
$\mu_k=v$ and $\sum_i\mu_ig_i=\sum_i\lambda_ig_i$ such that the face
$\{\eps_k=v\}$ contains no bad vertex.
\end{claim}
\begin{proof}
If there is no bad vertex, use Lemma~\ref{lem:slide}(i).

\emph{Suppose that $(0,0,1)$ is bad.}  The faces that do not contain this
vertex are $\{\eps_3=0\}$, $\{\eps_1=1\}$ and $\{\eps_2=1\}$.

If $\beta=3$, then $\alpha,\gamma\in\{1,2\}$, and
Lemma~\ref{lem:slide}(ii) gives $\mu_3=0$, $\mu_\alpha=1$ or
$\mu_\gamma=1$.  All three faces are permitted.

If $\beta\ne3$, name the indices so that $\gamma=3$; then
$\alpha,\beta\in\{1,2\}$.  Lemma~\ref{lem:slide}(i) gives $\mu_\beta=1$ or
$\mu_3=0$, which are permitted, or $\rho>0$ and $\mu_\alpha=0$.  In the
last case $\tau=\lambda_\alpha/\rho$, so by \eqref{eq:rel}
$\lambda_\alpha g_\alpha=\tau\,(g_\beta-\sigma g_3)$ and
\[
 \pi(0,0)=-\lambda_\alpha g_\alpha-\lambda_\beta g_\beta
 =-(\lambda_\beta+\tau)\,g_\beta+\tau\sigma\,g_3
 =-\mu_\beta\,g_\beta+(\lambda_3-\mu_3)\,g_3 .
\]
Here $\mu_\beta\in[0,1]$, and $\lambda_3-\mu_3=\tau\sigma\in[0,1]$ because
$0\le\mu_3\le\lambda_3$.  So each coordinate of $\pi(0,0)$ is a difference
of two numbers in $[0,1]$, that is, $\pi(0,0)\in Q$.  This contradicts the
assumption that $(0,0,1)$ is bad.

\emph{Suppose that $(1,1,0)$ is bad.}  The faces that do not contain this
vertex are $\{\eps_3=1\}$, $\{\eps_1=0\}$ and $\{\eps_2=0\}$.

If $\beta=3$, Lemma~\ref{lem:slide}(i) gives $\mu_3=1$, $\mu_\alpha=0$ or
$\mu_\gamma=0$, all permitted.

If $\beta\ne3$ and $\gamma=3$, Lemma~\ref{lem:slide}(ii) gives
$\mu_\beta=0$ or $\mu_3=1$, which are permitted, or $\rho>0$ and
$\mu_\alpha=1$.  In the last case $\tau=(1-\lambda_\alpha)/\rho$, so
$(1-\lambda_\alpha)g_\alpha=\tau\,(g_\beta-\sigma g_3)$ and
\[
 \pi(1,1)=(1-\lambda_\alpha)\,g_\alpha+(1-\lambda_\beta)\,g_\beta
 =(1-\mu_\beta)\,g_\beta-(\mu_3-\lambda_3)\,g_3 ,
\]
with $1-\mu_\beta\in[0,1]$ and $\mu_3-\lambda_3=\tau\sigma\in[0,1]$.  So
$\pi(1,1)\in Q$, a contradiction.
\end{proof}

\begin{proof}[Proof of Theorem~\ref{thm:upper}]
We show that $I(t)$ implies $I(t+1)$.  Take $k$, $v$ and $\mu$ from
Claim~\ref{cl:face}, and let $\{i,j\}=\{1,2,3\}\setminus\{k\}$.  Since
$\sum_l(\mu_l-\lambda_l)g_l=0$ and $\mu_k=v$, we have
$(v-\lambda_k)g_k=-(\mu_i-\lambda_i)g_i-(\mu_j-\lambda_j)g_j$, so for
every $\eps$ with $\eps_k=v$
\[
 P(\eps)=(\eps_i-\mu_i)\,g_i+(\eps_j-\mu_j)\,g_j .
\]
No such $\eps$ is bad, so every one of these four points that lies in $Q$
belongs to $R_{t+1}$.  This is $I(t+1)$, with generators $g_i,g_j$ and
coefficients $\mu_i,\mu_j$.  By induction $I(n)$ holds, and
Lemma~\ref{lem:nonempty} gives $R_n\ne\emptyset$.

\emph{Support.}  If $w_{t+1}=0$, the move $y_{t+1}=0$ does not move the
point, so $R_t\subseteq R_{t+1}$ along paths with $y_{t+1}=0$, and the old
parallelogram witnesses $I(t+1)$.  The case $w_{t+1}=1$ is the same with
$y_{t+1}=1$.  So the assignment can respect the support of $x$.

\emph{Algorithm.}  Store $g_1,g_2,\lambda_1,\lambda_2$ and, for each
vertex in $Q$, a pointer to a path that reaches it.  One step needs a
constant number of arithmetic operations: test the eight points, find the
relation \eqref{eq:rel}, compute $\tau$, and select the face.
\end{proof}

\begin{remark}
The proof uses the non-negativity of the weights twice: in the fact that
$\pi(1,0)$ and $\pi(0,1)$ lie in $Q$, and in Lemma~\ref{lem:middle}.  It
uses the dimension in Lemma~\ref{lem:middle}.
\end{remark}

\section{The lower bound}\label{sec:lower}

Let $k\ge0$ and $0\le\eta\le2^{-k}$.  The instance $\mathcal R_k(\eta)$
has $2k+1$ columns.  In the planar form all $w_j$ equal $1/2$, and the
vectors $h_j$ are, in order,
\[
 (1,\eta),\qquad\text{then for } i=k,k-1,\dots,1:\quad (1-2^{-i},\,1),\ \ (1,1).
\]

\begin{proof}[Proof of Theorem~\ref{thm:lower}]
We show that the optimal value of $\mathcal R_k(\eta)$ is
$T_k=1-2^{-(k+1)}$; for $\eta>0$ all weights are positive.  Each column
moves the point by $\pm h_j/2$, and the instance is invariant under
$p\mapsto-p$.  After the first column the point is $\pm(1/2,\eta/2)$; by
symmetry assume it is $(1/2,\eta/2)$.

Number the pairs $j=0,\dots,k-1$ in order, so that pair $j$ consists of
the vectors $(1-\eps_j,1)$ and $(1,1)$ with $\eps_j=2^{j-k}$, and put
\[
 u_j=\tfrac12+(2^j-1)\,2^{-(k+1)},\qquad S_j=(u_j,\eta/2).
\]
Then $S_0=(1/2,\eta/2)$.  Suppose that the point is $S_j$ before pair
$j$.  There are four ways to play the pair.
\begin{itemize}
\item First move $+$.  The first coordinate becomes
$u_j+(1-\eps_j)/2=T_k$.
\item Moves $-,-$.  The second coordinate becomes $\eta/2-1$, of absolute
value $1-\eta/2\ge T_k$.
\item Moves $-,+$.  The point becomes $(u_j+\eps_j/2,\eta/2)=S_{j+1}$.
\end{itemize}
Hence a path either attains sup norm at least $T_k$ inside some pair, or
it arrives at $S_k=(T_k,\eta/2)$.  So every path has a prefix of sup norm
at least $T_k$.  The path that plays $+$ and then $-,+$ in every pair
visits only the points $S_j$ and $(u_j-(1-\eps_j)/2,\,\eta/2-1/2)$, all of
sup norm at most $T_k$.
\end{proof}

The condition on $\eta$ cannot be dropped: $\mathcal R_7(1/100)$ has
optimal value $199/200=1-\eta/2$.

\begin{remark}
Since all $w_j$ are $1/2$, the upper bound $1$ for the instances
$\mathcal R_k$ is already contained in \citep{LundMagazinov2015}.  There
it is noted that the bound is attained in the max-norm plane, by two
vectors whose coordinates have mixed signs.  The family $\mathcal R_k$
shows that it is also approached by vectors with non-negative
coordinates, which are the only ones that two-machine instances produce.
We do not know whether the value $1$ is attained by a two-machine
instance; in exhaustive-evaluation searches the best values found with
$n=1,\dots,7$ columns are
$\tfrac12,\tfrac23,\tfrac34,\tfrac56,\tfrac78,\tfrac{11}{12},\tfrac{15}{16}$.
\end{remark}

\section{Online assignment}\label{sec:online}

\begin{proof}[Proof of Theorem~\ref{thm:online}]
Use the planar form and fix $c$.  Let $L$ be the set of points of
$[-c,c]^2$ from which the adversary has a strategy that forces the point
out of $[-c,c]^2$ after finitely many columns, and let
$W=[-c,c]^2\setminus L$.  If $p\in W$, then for every column $(h,w)$ at
least one of $p+(1-w)h$ and $p-wh$ lies in $W$; otherwise the adversary
presents this column and continues with the strategy of the point that
is reached, so $p\in L$.  We show that $W=\emptyset$.  Suppose that
$W\ne\emptyset$, and let
$M=\sup\{p_1:p\in W\}\le c$.  Fix an integer $K>4c$ and
$\delta=3^{-K}/4$, and choose $P\in W$ with first coordinate above
$M-\delta$.

For $i=0,\dots,K-1$ let $\delta_i=3^i\delta$ and present the column
$h=(4\delta_i,1)$, $w=1/2$; note that $4\delta_i\le1$.  Before step $i$ the point is in $W$ and its
first coordinate is above $M-\delta_i$.  The move $+$ raises the first
coordinate by $2\delta_i$, above $M$, so it leaves $W$.  Hence the move
$-$ leads to a point of $W$; its first coordinate is above
$M-3\delta_i=M-\delta_{i+1}$, and its second coordinate has dropped by
$1/2$.  After $K$ steps the second coordinate has dropped by $K/2>2c$, so
the point is outside $[-c,c]^2$, although it is in $W$.  This
contradiction shows $W=\emptyset$.  In particular the origin is in $L$.
\end{proof}

In the original variables the columns of this proof have
$d_{1j}=4\delta_i$, $d_{2j}=1$ and $x_{1j}=x_{2j}=1/2$, so all weights are
positive.

\section{Mixed signs, and open questions}\label{sec:open}

\begin{proposition}\label{prop:mixed}
Consider the planar walk of Section~\ref{sec:planar} with vectors
$h_j\in[-1,1]^2$ in place of $h_j\in[0,1]^2$.  For the six columns
$(h_j,w_j)$
\[
 \bigl((1,-1),\tfrac12\bigr),\
 \bigl((1,1),\tfrac47\bigr),\
 \bigl((1,-1),\tfrac17\bigr),\
 \bigl((\tfrac37,-1),\tfrac27\bigr),\
 \bigl((\tfrac9{14},-1),\tfrac5{14}\bigr),\
 \bigl((1,-1),\tfrac5{14}\bigr)
\]
every one of the $2^6$ paths leaves $Q$; the least possible maximal sup
norm is $15/14$.
\end{proposition}
\begin{proof}
By exhaustive evaluation in exact rational arithmetic; see the
repository.
\end{proof}

Such a column is not a column of a chairman instance, so
Proposition~\ref{prop:mixed} says nothing about $\kappa_m$.  It shows that
Theorem~\ref{thm:upper} is not a formal consequence of the symmetric
theorem of \citep{LundMagazinov2015}.

\begin{question}
What is the least $m$ with $\kappa_m>1$?  By Theorem~\ref{thm:upper} and
\citep{Muffatti2026} it satisfies $3\le m\le11$.  In particular, is
$\kappa_3=1$?
\end{question}

Searches with three and four rows and up to nine columns found values
above the two-row optimum for the same number of columns (for example
$0.9435$ with three rows and seven columns, against $15/16$), but none
above $1$.  These searches do not converge beyond eight columns and are
not evidence for $\kappa_3\le1$.

\begin{question}
Is the value $D$ attained by some two-machine instance?
\end{question}

\begin{question}
For two machines, can one keep $|\disc_1(t)+\disc_2(t)|\le D$ in addition
to $|\disc_i(t)|\le D$?  In the planar form this asks for a path in the
hexagon $\max(|p_1|,|p_2|,|p_1-p_2|)\le1$.  Searches found no
counterexample, but the invariant of Section~\ref{sec:upper} does not
carry over: its closure fails for the hexagon in random tests.
\end{question}

\section*{Reproducibility and disclosure}

The accompanying repository (directory \texttt{chairman-small-rows})
contains, in exact rational arithmetic and with the Python standard
library only: an implementation of the proof of Theorem~\ref{thm:upper}
as an algorithm that asserts Claims~\ref{cl:bad} and~\ref{cl:face},
Lemmas~\ref{lem:middle} and~\ref{lem:slide} and the invariant at every
step, run on $40{,}048$ instances and compared with exhaustive search on
$2{,}665$ small ones; an exhaustive check of Theorem~\ref{thm:lower} for
$k\le9$; and the check of Proposition~\ref{prop:mixed}.  The search
results quoted in Sections~\ref{sec:lower} and~\ref{sec:open} come from
simulated annealing and are not certified optima.

Generative AI (Claude, Anthropic) was used extensively in this work: for
the computational exploration, for finding and writing the proofs,
including the invariant of Section~\ref{sec:upper}, for the verification
code, and for an adversarial review of the proof.  The AI systems are not
authors.  The author assumes responsibility for independent verification,
priority, attribution, and the final text.  At the date of this draft the
proofs have been checked by machine tests and by AI review, but not yet
by an independent human reader or a proof assistant.

\bibliographystyle{abbrvnat}
\bibliography{references}

\end{document}
